\documentclass[11pt]{article}
\usepackage[letterpaper, margin=0.85in, top=0.9in, bottom=0.95in, headsep=16pt]{geometry}
\usepackage{libertinus}
\usepackage[scale=0.82]{sourcecodepro}
\usepackage[table]{xcolor}
\usepackage{fvextra, booktabs, longtable, enumitem, titlesec, fancyhdr, microtype, amsmath, array}
\usepackage[most]{tcolorbox}
\usepackage{needspace}
\usepackage[hidelinks]{hyperref}

% J.T.S. palette
\definecolor{ink}{HTML}{1E2B24}
\definecolor{muted}{HTML}{56645B}
\definecolor{accent}{HTML}{2B5E8C}
\definecolor{accentsoft}{HTML}{DCE6EF}
\definecolor{moss}{HTML}{4F6E4C}
\definecolor{mosssoft}{HTML}{E1EADC}
\definecolor{sand}{HTML}{CDBF9E}
\definecolor{surfalt}{HTML}{E7EDE3}
\definecolor{panel}{HTML}{F7F9F5}
\definecolor{rulec}{HTML}{D2DACD}
\definecolor{warn}{HTML}{8A6410}
\color{ink}

\newcommand\Rl[1]{\textcolor{accent}{\textbf{#1}}}
\newcommand\Ms[1]{\textcolor{moss}{\textbf{#1}}}
\newcommand\Dm[1]{\textcolor{muted}{#1}}

\fvset{fontsize=\footnotesize, breaklines, breakanywhere, breakindent=2.5em,
  breaksymbolleft={\tiny\textcolor{sand}{\ensuremath{\hookrightarrow}}}}

\newtcolorbox{termbox}{breakable, enhanced, colback=panel, colframe=rulec, boxrule=0.4pt,
  borderline west={2.2pt}{0pt}{sand}, arc=2pt, left=7pt, right=5pt, top=4pt, bottom=4pt,
  before skip=6pt, after skip=10pt}
\newtcolorbox{thmbox}[1]{breakable, enhanced, colback=accentsoft!45, colframe=accent!55, boxrule=0.5pt,
  arc=2pt, left=8pt, right=8pt, top=5pt, bottom=5pt, before skip=8pt, after skip=10pt,
  title={\sffamily\small\bfseries #1}, coltitle=white, colbacktitle=accent, fonttitle=\sffamily}
\newtcolorbox{keybox}{enhanced, colback=mosssoft!55, colframe=moss!45, boxrule=0.5pt, arc=2pt,
  left=9pt, right=9pt, top=6pt, bottom=6pt, before skip=8pt, after skip=8pt}

\newcommand\cmdtag[1]{\colorbox{accentsoft}{\textcolor{accent}{\ttfamily\bfseries #1}}}
\newcommand\proved{\textcolor{moss}{\textbf{proved}}}
\newcommand\onestep{\textcolor{moss}{\textbf{one step}}}
\newcommand\stepok{\textcolor{accent}{\textbf{step}}}

\titleformat{\section}{\sffamily\Large\bfseries\color{ink}}{}{0pt}{}[\vspace{2pt}{\color{sand}\titlerule[0.8pt]}]
\titlespacing*{\section}{0pt}{20pt}{9pt}
\titleformat{\subsection}{\sffamily\large\bfseries\color{ink}}{}{0pt}{}
\titlespacing*{\subsection}{0pt}{14pt}{5pt}

\pagestyle{fancy}
\fancyhf{}
\renewcommand\headrulewidth{0pt}
\fancyhead[L]{\sffamily\small\color{muted}Collins' projection in pvs\_cad}
\fancyhead[R]{\sffamily\small\color{muted}\thepage}

\setlength\parindent{0pt}
\setlength\emergencystretch{3em}
\setlength\parskip{5pt}
\newcommand\code[1]{\mbox{\texttt{#1}}}
% no hyphenation (and no break at a hyphen) inside typewriter identifiers
\makeatletter
\AddToHook{selectfont}{\edef\jts@f{\f@family}\edef\jts@t{\ttdefault}%
  \ifx\jts@f\jts@t\hyphenchar\font=-1\relax\fi}
\makeatother
\renewcommand\arraystretch{1.25}

\begin{document}
\thispagestyle{empty}
\begin{flushleft}
{\sffamily\small\bfseries\color{moss}\MakeUppercase{pvs\_cad \quad·\quad PVS 8.1 \quad·\quad 3 October 2026}}\par\vspace{8pt}
{\fontsize{27}{33}\selectfont\bfseries Collins' projection, verified}\par\vspace{10pt}
{\large\color{muted} The cylindrical algebraic decomposition of pvs\_cad, built from Collins' projection: the
theorem that the projections form a CAD, the decision procedure and the quantifier elimination that read truth
values off it, what they cost on the benchmarks, and where the limits are.}
\end{flushleft}
\vspace{4pt}

\begin{keybox}
\textbf{In one line.} For every family of polynomials in any number $n$ of variables, the cells built from Collins'
projection are \textbf{proved} to form a cylindrical algebraic decomposition adapted to the family --- with no
certificate checked at run time --- and every cell is defined without quantifiers by sign conditions on polynomials
with rational coefficients, so it is semi-algebraic in the sense of Basu--Pollack--Roy~\cite{bpr}. An executable function
returns that CAD as data: at least one record, each with an exact sample point, for every nonempty cell of
$\mathbb{R}^n$, and none for any other. The decision
procedure built on it (\texttt{decide8}) is \textbf{proved} sound and complete, and its quantifier elimination
(\texttt{qe8}) is \textbf{proved} to return a quantifier-free equivalent of every prenex formula with at least one
quantifier and one or more free variables; \texttt{qelim}, built on both, does the same for every first-order
formula. \texttt{(cad)} and \texttt{(cad-qe)} use them in a theory that imports \texttt{pvs\_cad}.
\end{keybox}

\section{What is proved}

\begin{thmbox}{col\_stack.collins\_stack \quad (delineability from Collins' projection)}
\begin{Verbatim}[fontsize=\small]
collins_stack: THEOREM conn?(C) AND C(p0) AND colh?(G, C) IMPLIES delin_cl?(G, C)
\end{Verbatim}
\end{thmbox}
If $C$ is connected and nonempty and Collins' N-reducta projection of the family $G$ is sign-invariant on $C$
(\code{colh?}: the leading coefficients of the kept reducta, the principal subresultant coefficients
$\mathrm{psc}_j(f, f')$, and $\mathrm{psc}_j(f, g)$ for pairs each have one sign on $C$), then $G$ is
\emph{delineable} over $C$ in the sign-invariant sense (\code{delin\_cl?}; root multiplicities are not tracked):
over each point of $C$, the real roots of the members not identically zero on its fibre are the values of a
constant number of continuous, strictly ordered functions on $C$, and every member has one sign on each section and band
between them. This is Collins' delineability theorem, the basis of Collins' original algorithm~\cite{collins} (the smaller
projections of McCallum~\cite{mccallum}, Brown~\cite{brown} and Lazard~\cite{lazard,mpp} rest on different theorems); here it is proved from the definitions
(subresultant theorem, continuity of the roots of a polynomial in its coefficients, connectedness).

\begin{thmbox}{col\_run.col\_cad\_run \quad (the projections form a CAD)}
\begin{Verbatim}[fontsize=\small]
col_cad_run: THEOREM k >= 1 IMPLIES
      (FORALL p: length(p) = k + 1 IMPLIES
         EXISTS s, A: ... ccell(k, F)(s, A)(p))                     % cover
  AND (... ccell(k, F)(s1, A1)(p) AND ccell(k, F)(s2, A2)(p)
         IMPLIES s1 = s2 AND A1 = A2)                              % disjoint
  AND (... conn?(ccell(k, F)(s, A))
         AND fod?(1 + length(A))(ccell(k, F)(s, A)))                % connected, definable
  AND (FORALL s: member(s, csects(k, F))
         IMPLIES cad_over?(ctw(F, k), scell(s)))                    % delineable
  AND (... the graph of every root function is definable ...)
  AND (... ccell(k, F)(s, A)(p1) AND ccell(k, F)(s, A)(p2)
         IMPLIES svec(F, p1) = svec(F, p2))                         % F sign-invariant
\end{Verbatim}
\end{thmbox}
For every family $F$ in $k+1$ variables: project $k$ times, $P_k = F$ and $P_{j-1}$ a basis of
$\mathrm{proj}(P_j)$, meant to be squarefree, that a run-time product check confirms factors every member
($\mathrm{proj}(P_j)$ itself when the check fails; \code{fsplit}) (\code{ctw(F, k)} is the tower $P_1, \dots, P_k$), cut the line of the outermost variable at the real roots of the
$k$-th projection $P_0$ (\code{pbot}, \code{csects}), and build the stacks. The cells of $\mathbb{R}^{k+1}$ cover it
and are disjoint; the cells of every level are connected and first-order definable; every level is delineable over
every cell below it, and the graphs of its root functions are definable; and every polynomial of $F$ has one sign
on every cell of $\mathbb{R}^{k+1}$. Nothing is checked at run time: the statement has no hypothesis but
$k \ge 1$, and \code{col\_cells} adds $k = 0$, one variable, where the cells are the sectors of the roots of $F$.
\code{col\_verified} adds the data: when the walk's certificate holds (two or more variables), the run returns at
least one record for every nonempty cell of $\mathbb{R}^{k+1}$ and none for any other, each with a sample point in
its cell and $F$'s signs there; \code{col\_found} raises the run's effort until the certificate holds, and
\code{col\_found\_cad} says it always returns the CAD as data, in every number of variables, one included. These
are the conditions of Basu--Pollack--Roy's Defs.~5.1 and~5.5 for a CAD adapted to $F$, with the partition stated for
$\mathbb{R}^{k+1}$ (at lower levels it follows from the stack construction, \code{tcell\_cover} and
\code{tcell\_disj}, and is not restated) and its finiteness given by the records. First-order definable sets
(formulas with rational coefficients, no parameters) are exactly the quantifier-free definable ones
(\code{fod\_qfd}, Tarski--Seidenberg, proved for every first-order formula), so every cell is semi-algebraic in
the sense of Basu--Pollack--Roy; indeed it is defined by a Boolean combination of sign conditions on polynomials
with rational coefficients (\code{col\_cells\_sa}), where their definition allows real coefficients.
\code{decb\_run.decb\_cad\_run} states the same for the cells of \texttt{decide8}'s run, which cuts the line of
the outermost variable at the roots of a smaller bottom-step projection of $P_1$ (\code{projb}) when every member
is nonzero at the rational sample of every sector that is not a root (\code{nzb?}), and at the roots of $P_0$
otherwise; \code{decb\_cad\_input} gives the input family one sign vector on every cell when \texttt{decide8} runs
over its checked basis.

{\small
\begin{longtable}{@{}>{\ttfamily}p{0.22\textwidth} p{0.72\textwidth}@{}}
\toprule
\textsf{\small\bfseries Theorem} & \textsf{\small\bfseries What it says}\\
\midrule
ctw\_cad & over any connected set on which the $k$-th projection is sign-invariant, the tower of projections is a CAD\\
col\_cad\_run & the projections form a CAD of $\mathbb{R}^{k+1}$ adapted to $F$, for every $F$ and $k\ge1$\\
col\_cells & the same for every $k$, one variable included\\
col\_verified & under the walk's certificate, the run's records name that CAD's cells of $\mathbb{R}^{k+1}$, at least one per nonempty cell, with a sample and $F$'s signs in each\\
col\_found\_cad & \texttt{col\_found}, with the search for the run's effort built in, always returns the CAD as data\\
col\_found\_decides & its records decide every sentence over $F$ in the same variables, in the same order\\
fod\_qfd & a set is first-order definable exactly when it is quantifier-free definable (Tarski--Seidenberg)\\
col\_cells\_sa & every cell at every level, and every root graph, is quantifier-free definable over $\mathbb{Q}$, hence semi-algebraic\\
decc\_correct & the decision over the Collins tower, when certified, is the truth of the sentence\\
decc\_complete & for at least two quantifiers some run is certified\\
projb\_delin & over a connected set where the members of the smaller bottom-step projection are nonzero, the lowest family is delineable\\
decb\_cad\_run & the cells of \texttt{decide8}'s run, with the smaller bottom step, form a CAD adapted to its family, for every $F$ and $k\ge1$\\
decb\_correct & \texttt{decide8}'s run, when certified, is the truth of the sentence; \code{decb\_complete}: some run is certified\\
decide8\_decides & \texttt{decide8} decides every closed prenex formula: its value is the truth (sound and complete)\\
qe8\_complete & quantifier elimination by the Collins run: a quantifier-free equivalent of every prenex formula with at least one quantifier and $m \ge 1$ free variables\\
qelim\_ok & \texttt{qelim} (with \texttt{qelim\_qf}): a quantifier-free equivalent of every first-order formula, quantifiers anywhere, any $m \ge 0$\\
pstower\_cad & over any connected set on which the projection of its lowest family is sign-invariant, a tower whose families \emph{contain} the projections of the next is a CAD (used for derivative-closed free levels)\\
detf\_det & the determinant computed with each minor once ($n\,2^n$ products) is the Laplace determinant ($n!$)\\
bdet\_det & the determinant by fraction-free elimination (Bareiss, $O(n^3)$), with each exact division checked, is the Laplace determinant (in normal form)\\
\bottomrule
\end{longtable}}

\textbf{What you trust.} The PVS kernel; its ground evaluator, which runs the procedures on the concrete input
inside the proof (proof by reflection), and the other evaluations the strategies use (normal forms, witness points,
NASALib's interval enclosures); NASALib's saved proofs of the theorems the library imports (from
\texttt{reals}, \texttt{Sturm}, \texttt{Tarski}, \texttt{structures}, \texttt{analysis}, \texttt{complex},
\texttt{matrices}, \texttt{mult\_poly}, \texttt{interval\_arith}, \texttt{trig} and \texttt{lnexp}, and through them
from other NASALib libraries such as \texttt{ints} and \texttt{power}), which were not replayed here; and the reading
of the statements (\code{fsem}, \code{gsem}, \code{svec}, \code{delin\_cl?}, \code{cad\_over?} say what is
intended). The procedures and the strategy code are not trusted; the procedures' correctness is the theorems.

\textbf{Size.} The Collins work (the theories listed under ``Where it is'', with two example theories) is 76
theories with 1,014 proved formulas (theorems, lemmas and type-correctness conditions) in 4,799 lines of PVS. The
whole library has 375 theories and 5,628 proved formulas, all replayed by \texttt{proveit} from a clean copy
(5 October 2026, 58 minutes).

\section{Using it}

{\raggedright In a theory, write \texttt{IMPORTING pvs\_cad} (\texttt{cad/pvs\_cad.pvs} imports
\texttt{cad\_decide8}, \texttt{qe8}, \texttt{gform\_ok}, \texttt{mpoly\_embed} and \texttt{cad\_endgame}).
\cmdtag{(cad)} and \cmdtag{(cad-direct)} then decide a closed formula. A prenex one (one real per quantifier;
grouped binders are decided through the nested form) is decided by \texttt{decide8}: the Collins run for two or
more quantifiers, the sectors of the line for one; the proof instantiates \texttt{decide8\_correct}. A formula of
any other shape (quantifiers inside \texttt{NOT}, \texttt{AND}, \texttt{OR}, \texttt{IMPLIES} or \texttt{IFF};
binders over \texttt{posreal}, \texttt{nnreal}, \texttt{negreal}, \texttt{npreal} or \texttt{nzreal}) is decided by
\texttt{decide\_g}, which eliminates its quantifiers from the inside out with \texttt{qe8} and \texttt{decide8}; the
proof instantiates \texttt{decide\_g\_ok}. \cmdtag{(cad :cad-only? t)} answers only through these decisions; with
one quantifier that is the sectors of the line, not the Collins run. \cmdtag{(cad-qe FNUM)} eliminates the
quantifiers of a formula with free variables by \texttt{qe8} (\texttt{qe8\_correct}), or, for a formula of any
other shape, by \texttt{qe\_g} (\texttt{qe\_g\_ok}). \texttt{(cad)} also has shortcuts that \texttt{(cad-direct)}
and \texttt{(cad :cad-only? t)} skip: a witness search, which proves a goal whose quantifiers are all existential
(or refutes a hypothesis whose quantifiers are all universal) by one point; with two or more quantifiers, a constant
for one variable (existential in a goal, universal in a refuted hypothesis) that settles the matrix outright; and, with three or more, when the full decision
does not answer within 10 s, one variable fixed to a constant and the smaller sentence decided. With
\texttt{cad\_decide5} imported instead of \texttt{cad\_decide8}, the decision is the read-closure \texttt{decide5}
(\texttt{decide7} under \texttt{:cad-only?}).\par}
\begin{termbox}
\begin{Verbatim}
c8_l3_quad: LEMMA FORALL (c: real): EXISTS (x1: real): EXISTS (x2: real): EXISTS (x3: real):
  (x1^2 <= 1 AND x2^2 <= 1 AND x3^2 <= 1 AND x1 + x2 + x3 >= c) OR c > 3
Rule? (cad-direct)
Q.E.D.                      [about 5 s; the read-closure decision did not finish in 600 s]
\end{Verbatim}
\end{termbox}
{\small\color{muted} \texttt{c8\_l3\_quad} is in \texttt{cad\_decide8\_ex}, which imports \texttt{cad\_decide8};
the transcript is abridged (the step's messages are left out).}

\section{What it costs}

\textbf{The decision alone} (\texttt{decc\_o}, the Collins run, against \texttt{decn\_o}, the read-closure run,
both at $u = 0$; \code{(cad-mx)} measurements on 1 October 2026, wall clock on one machine; ``$>$'' marks a run
stopped at that time):

{\small
\begin{longtable}{@{}>{\ttfamily}p{0.15\textwidth} c p{0.38\textwidth} r r@{}}
\toprule
\textsf{\small\bfseries Formula} & \textsf{\small\bfseries quantifiers} & \textsf{\small\bfseries shape} & \textsf{\small\bfseries Collins} & \textsf{\small\bfseries read closure}\\
\midrule
t\_circle & 2 & $\forall x\,\exists y.\; x^2+y^2=1 \lor x^2>1$ & 0.1 s & 0.0 s\\
t\_sphere & 3 & the same in three variables & 0.1 s & 0.1 s\\
e4\_above & 4 & $\forall x\,y\,z\;\exists w.\; w>x \land w>y \land w>z$ & 2.1 s & 0.9 s\\
e4\_between & 4 & the same with $w<z$ (false) & 0.2 s & 0.4 s\\
k2\_lin & 3 & box $[-1,1]^2$ against $x_1+x_2 \ge c$, $c$ universal & 0.1 s & 0.2 s\\
l2\_quad & 3 & the same with $x_i^2 \le 1$ & 0.1 s & 0.8 s\\
k3\_lin & 4 & box $[-1,1]^3$ against $x_1+x_2+x_3 \ge c$ & 0.5 s & 11.5 s\\
l3\_quad & 4 & the same with $x_i^2 \le 1$ & 0.7 s & $>$600 s\\
bath\_01 & 3 & ball and cylinder (Bath CAD bank~\cite{bath}) & 0.2 s & $>$90 s\\
bath\_08 & 3 & Random B, Gr\"obner form (false) & 0.8 s & 1.1 s\\
bath\_02\textrm{--}05, 07, 09, 10, 12 & 3--4 & the other Bath problems & $>$90 s & $>$90 s\\
\bottomrule
\end{longtable}}

\textbf{Every proof of the example theories}, replayed through \texttt{(cad)}, \texttt{(cad-direct)} or
\texttt{(cad-qe)}, including the witness search and the reconstruction of the proof: as they are, importing
\texttt{cad\_decide5} (\texttt{qe\_all} for \texttt{qe\_cad\_ex}), and with \texttt{cad\_decide8} (\texttt{qe8} for
\texttt{qe\_cad\_ex}) added to their imports. Measured on 1 October 2026 with the strategies of that day, before
the changes for formulas as people write them:

{\small
\begin{longtable}{@{}>{\ttfamily}p{0.2\textwidth} r r r@{}}
\toprule
\textsf{\small\bfseries Theory} & \textsf{\small\bfseries proofs} & \textsf{\small\bfseries cad\_decide5 / qe\_all} & \textsf{\small\bfseries cad\_decide8 / qe8}\\
\midrule
cad\_examples\textrm{, }cad\_examples2\textrm{--4} & 28 & 69.1 s & 90.2 s\\
cad\_endgame\_ex & 18 & 31.0 s & 43.0 s\\
cad\_star\_ex & 15 & 39.9 s & 49.5 s\\
cad\_bath & 1 & 9.5 s & 10.2 s\\
qe\_cad\_ex & 6 & 42.0 s & 51.1 s\\
\midrule
all & 68 & 191.5 s, 68 proved & 244.0 s, 68 proved\\
\bottomrule
\end{longtable}}

In that run every proof went through either way, and again in a re-run after the strategy fixes of this release (68 of 68 with
each import, 1 October 2026). The whole \texttt{(cad -1)} proof of \texttt{cad\_bath.bath\_08\_false} (Bath 08,
``Random B, Gr\"obner form'', false) is one of them: 9.5\,s as it is and 10.2\,s with \texttt{cad\_decide8}. On these
small examples, in that first run, a proof took 0.77 s longer on average with \texttt{cad\_decide8} or \texttt{qe8}
imported: 0.1--1.3 s for a \texttt{(cad)} or \texttt{(cad-direct)} proof, 0.5--2.5 s for a \texttt{(cad-qe)} proof (in the
re-run after the strategy fixes, 0.59 s on average: 166.7 s against 206.9 s for the 68 proofs). Part of that is
the larger imported context, not the Collins run: a sentence with one quantifier runs the same decision,
\texttt{decq2}, under either import, and the 16 one-quantifier proofs of \texttt{cad\_endgame\_ex} are 0.2--1.1 s
slower (0.64 s on average). On some problems with several quantifiers the Collins run is faster by more than an
order of magnitude (\texttt{k3\_lin}: 23 times; \texttt{bath\_01} and \texttt{l3\_quad}: more than 450 times, and
the whole proof of \texttt{l3\_quad} through \texttt{(cad-direct)} takes about 5 s, against more than 600 s for the
read-closure decision alone); on others it is slower (\texttt{e4\_above}: 2.1 s against 0.9 s).


\section{How it works}
\begin{enumerate}[leftmargin=1.3em, itemsep=3pt]
\item \textbf{The projection} (\code{projn}): for each polynomial of the family, the leading coefficients of its
kept reducta (its coefficients from the top down to the first that is a nonzero constant) and the principal
subresultant coefficients $\mathrm{psc}_j(f, f')$ of each kept reductum $f$; for every pair of polynomials,
$\mathrm{psc}_j(f, g)$ for their kept reducta $f$ and $g$; normalized, without constants and repetitions. The
psc are determinants of Sylvester-type matrices, computed by fraction-free elimination (\code{bdet}). Each
projection is then replaced by a basis of it, meant to be squarefree (\code{fsplit}), when a product check confirms
that the basis factors every member; the basis has the same roots, so the CAD is the same (\code{fsplit\_cinv}).
\item \textbf{Delineability} (\code{collins\_stack}): over a connected set where the projection is sign-invariant,
the number of distinct roots is constant (equal psc patterns of $(f, f')$ give equal numbers of distinct roots),
the roots move continuously (lower semicontinuity of the roots in the coefficients), one root stays in each small
disc, and common roots of two members stay common (the subresultant $S_\kappa$); so the root functions are
continuous and never cross.
\item \textbf{The tower}: projecting $k$ times gives $P_{k-1}, \dots, P_0$ below $P_k = F$; over each sector of
the roots of $P_0$ the stacks of $P_1, \dots, P_k$ form the CAD (\code{ctw\_cad}), because each level's
projection is sign-invariant on the cells of the level below. \texttt{decide8} cuts the line at the roots of a
smaller bottom step instead (\code{projb}: each member's leading coefficient and, for each member with its derivative and each pair of
members, only the first psc that is not identically zero), which is
enough where its members are nonzero (\code{projb\_delin}, \code{decb\_cad\_run}).
\item \textbf{Truth values} come from the verified walk of the read-closure engine, over exact algebraic sample
points: it passes the sections and bands of each level and folds the quantifiers. Its certificate only checks the answers of
the sign oracle at the samples; delineability no longer depends on it.
\end{enumerate}

\section{Limits}
\begin{itemize}[leftmargin=1.3em, itemsep=3pt]
\item \textbf{The size of Collins' projection.} It contains psc of every pair of polynomials and of every kept
reductum. On \texttt{bath\_04} (one polynomial in three variables, of degree 4 in its main variable) the first
projection has 11 polynomials and the second 569. Two remedies are in place: every family is replaced by a checked
basis (square-free factors), and \texttt{decide8} cuts the line at the roots of the smaller bottom step; in a
prototype count \texttt{bath\_04}'s bottom family drops from 569 polynomials to 35. The time then goes to the lifting,
and \texttt{bath\_04}, \texttt{05} and \texttt{07} are still not decided in 300 s (3 October 2026). Not done: a
smaller projection at every level (McCallum's~\cite{mccallum} or Brown's~\cite{brown}; \texttt{cad\_proj} defines a
McCallum-shaped operator, not proved delineating; the read-closure engine starts from it, under its run-time
certificate), Lazard's~\cite{lazard,mpp}, and equational constraints. Measured with
\texttt{decide5} (\texttt{cad\_decide5}), September 2026: \texttt{(cad)} settles Bath 01, 03 and 04 by a model
search, and Bath 02 --- and 05 and 07, refuted as hypotheses --- by fixing one variable to a constant and deciding
the smaller sentence, after the full decision does not answer within its 10 s budget. With \texttt{cad\_decide8}
imported (\texttt{IMPORTING pvs\_cad}, 1 October 2026) the same six are settled the same way: Bath 01, 03 and 04 in
0.5--0.6\,s, Bath 02, 05 and 07 in about 12\,s each (whole proofs; the smaller sentences are decided by
\texttt{decide8}).
\item \textbf{Universal sentences with many polynomials} (the Joukowsky problems \texttt{bath\_09}, \texttt{10}
and \texttt{12}) get no help from a witness search, and their decision does not finish in 90 s with either
engine.
\item \textbf{Quantifier elimination with two or more free variables} first answers by sign conditions on the
projection polynomials of the plain tower; when they do not tell the true cells from the false ones (two cells
with the same signs and different truth values), \texttt{qe8} runs the Collins tower whose free families are
closed under the derivative (\texttt{ctwd}), where by Thom's lemma they always do. The answer is always
quantifier-free; the derivative-closed tower is larger, so such problems cost more.
\item \textbf{Determinants} are computed by fraction-free elimination (Bareiss): $O(n^3)$ polynomial operations
instead of $n\,2^n$. Each exact division is done by unproved code and kept only when the quotient times the
divisor has the normal form of the dividend; when a check fails, the determinant is computed with every minor
once ($n\,2^n$ products), so the result is always the determinant. The proof (\code{bdet\_det}) runs through
Chi\`o's identity, from column operations, and the absence of zero divisors among polynomials
(\code{mpoly\_dom}).
\item \textbf{Where it is.} The theories are in \texttt{cad/}: the \texttt{col\_*}, \texttt{decc\_*},
\texttt{cad\_decide8*}, \texttt{qe1c\_*}, \texttt{qe2cc\_*}, \texttt{qe2cd\_*}, \texttt{ctwd*}, \texttt{twrun\_*},
\texttt{qe8*} and \texttt{qelim\_*} files and \texttt{cad\_sa}, with the algebra they rest on (\texttt{sres\_*},
\texttt{cpoly\_*}, \texttt{croot\_*}, \texttt{ev\_near}, \texttt{rdet\_*}, \texttt{det\_fast*}, \texttt{bareiss*}, \texttt{mpoly\_dom},
\texttt{cad\_projn*}), and the smaller bottom step and the checked basis (\texttt{cad\_projb*}, \texttt{col\_towerb},
\texttt{decb\_*}, \texttt{mpoly\_gcd\_def}, \texttt{mpoly\_cert*}, \texttt{fsplit\_top*}).
\texttt{cad/top.pvs} lists and describes every theory of the library.
\end{itemize}

\begin{thebibliography}{9}\footnotesize
\setlength\itemsep{1pt}
\bibitem{collins} G.~E. Collins. Quantifier elimination for real closed fields by cylindrical algebraic
  decomposition. In \emph{Automata Theory and Formal Languages}, LNCS 33, pages 134--183. Springer, 1975.
\bibitem{bpr} S.~Basu, R.~Pollack, M.-F.~Roy. \emph{Algorithms in Real Algebraic Geometry}, 2nd ed. Springer, 2006.
\bibitem{mccallum} S.~McCallum. An improved projection operation for cylindrical algebraic decomposition. In
  B.~F. Caviness, J.~R. Johnson (eds.), \emph{Quantifier Elimination and Cylindrical Algebraic Decomposition}, pages
  242--268. Springer, 1998.
\bibitem{brown} C.~W. Brown. Improved projection for cylindrical algebraic decomposition. \emph{Journal of Symbolic
  Computation} 32(5):447--465, 2001.
\bibitem{lazard} D.~Lazard. An improved projection for cylindrical algebraic decomposition. In C.~L. Bajaj (ed.),
  \emph{Algebraic Geometry and its Applications}, pages 467--476. Springer, 1994.
\bibitem{mpp} S.~McCallum, A.~Parusi\'nski, L.~Paunescu. Validity proof of Lazard's method for CAD construction.
  \emph{Journal of Symbolic Computation} 92:52--69, 2019.
\bibitem{bath} D.~J. Wilson, R.~J. Bradford, J.~H. Davenport. A repository for CAD examples. \emph{ACM Communications
  in Computer Algebra} 46(3/4):67--69, 2012.
\end{thebibliography}


\end{document}
